\cotefile{158}
\scanpage{20}
\mspage{1}

\uline{Proposition 1} Soit $X\in\mathrm{Ob}\,\mathrm{Cat}$ \deleted{\unsure{et $0$-co}} une catégorie pseudo-cofiltrante ($\Longleftrightarrow X^\wedge$ est localement irréductible), ayant un plus petit objet $e$\inserted{(donc $X$ est \unsure{$0$-conn.|$0$-connexe}).}

Soit $X_0$ la sous-catégorie pleine réduite à $e$, $M=\mathrm{End}_X(e)=\mathrm{End}_{X_0}(e)$ le monoïde des endomorphismes de $e$, $G=\pi_1(M)=\pi_1(X_0,e)$ le groupe enveloppant de $M$. Alors

a) L'hom canonique $\pi_1(X_0,e)\longrightarrow\pi_1(X,e)$ est surjectif, \deleted{et \unsure{son|pour}}

b) Pour toute double-flèche de source $e$
\diagram{\begin{tikzcd}[ampersand replacement=\&]e\arrow[r,shift left=.55ex,"u"]\arrow[r,shift right=.55ex,"v"']\&x\end{tikzcd}}
dans $X$, il existe $u',v'\in M$ tels que
\[
uu'=vv',
\]
et l'élément
\[
\unsure{[v]^{-}[u]|[v]^{-1}[u]|[v][u]}\ \text{de}\ \pi_1(X,e)
\]
(image canonique du chemin \diagram{\begin{tikzcd}[ampersand replacement=\&]e\arrow[r,"u"]\&x\&e\arrow[l,"v"']\end{tikzcd}})
est égale à l'image dans $\pi_1(X,e)$ de l'élément $[v,u]\overset{\mathrm{def}}{=}v'u'^{-1}\in G=\pi_1(X_0,e)$.

De plus, \inserted{si $e$ est aussi objet maximal de $X$,} cet élément \inserted{$[v,u]$} de $G$ ne dépend pas du choix du couple $(u',v')$ \deleted{\unsure{d'}}d'éléments de $M$ satisfaisant $uu'=vv'$. Plus précisément, si $\alpha:x\longrightarrow e$, alors $\deleted{\unsure{w u=v=}}u_1=\alpha u$ et $v_1=\alpha v$ sont des éléments $\deleted{[v,u]=}$ de $M$, et on a $u_1u'=v_1\unsure{u'|v'}$ d'où
\[
[v,u]=v_1^{-1}u_1^{\deleted{\unsure{-1}}}.
\]

c) Pour toute double-flèche \diagram{\begin{tikzcd}[ampersand replacement=\&]e\arrow[r,shift left=.5ex,"u"]\arrow[r,shift right=.5ex,"v"']\&x\end{tikzcd}}, dans $X$,

% archive sheet number: 19
\scanpage{21}
\mspage{\textcircled{2}}

choisissons $u',v'$ \inserted{\unsure{dans|de} $M$} tels que $uu'=vv'$, \inserted{d'où $[u,v]\in G$.} Alors le noyau de l'hom surjectif $G=\pi_1(X_0,e)\longrightarrow\pi_1(X,e)$ est le \inserted{ss-}groupe invariant de $G$ engendré par les éléments suivants, correspondants aux diagrammes (pas nécessairement commutatifs)
\diagram{\begin{tikzcd}[ampersand replacement=\&]x\arrow[r,"f"]\&y\arrow[r,"g"]\&z\\[-1ex]\&e\arrow[ul,"u"]\arrow[u,"v"]\arrow[ur,"w"']\&\end{tikzcd}}
\margin{\uline{NB} S'il existe $\alpha:z\longrightarrow e$, on voit aussitôt que cette relation est triviale. Donc si $e$ est aussi plus grand objet de $X$, alors $\pi_{\unsure{1|0}}(X_0,e)\xrightarrow{\sim}G\xrightarrow{\sim}\pi_1(X,e)$ iso !}
\[
[w,gv][v,fu][w,gfu]^{-1}.
\]

Soient $M$ un monoïde \inserted{\fbox{pseudocofiltrant}$^{*}$}, \inserted{$Y=B_M$,} $\varphi\in M$, considérons la relation d'équivalence \inserted{$R=R_\varphi$} dans $M$
% slanted
\margin{i.e. pour $u,v\in M$, $\exists\,u',v'\in M$ \unsure{avec} $uu'=vv'$}
\[
u\overset{R}{\sim}v\quad\Longleftrightarrow\quad\exists\,p,q\in\mathbb N\qquad\varphi^p u=\varphi^q v.
\]
Cette relation d'équivalence est stable par translations à droite, donc
\[
F\overset{\mathrm{def}}{=}M/\unsure{\mathcal R|R}
\]
est un $M$-ens à droite, i.e un objet de $\deleted{\unsure{X}}Y^\wedge$. On pose $X=Y_{/F}$, donc
% letterhead
\par\smallskip
{\footnotesize\noindent\begin{minipage}[t]{.46\linewidth}\raggedright
\deleted{MINISTERE DE L'EDUCATION NATIONALE}\\ \deleted{ACADEMIE D'AIX-MARSEILLE}\\ \rule{2.6em}{.3pt}\\ \deleted{Inspection Académique}\\ \deleted{de Vaucluse}\\ \rule{2.6em}{.3pt}\\ Tél. : 90 82 99 76\\ \deleted{Référence : Secrétariat n\textsuperscript{o}}
\end{minipage}\hfill%
\begin{minipage}[t]{.51\linewidth}\raggedright
\deleted{AVIGNON, le}\\ \deleted{L'Inspecteur d'académie,}\\ \deleted{directeur des services départementaux}\\ \deleted{de l'Education Nationale de Vaucluse,}\\ \deleted{à M}
\end{minipage}\par}

% archive sheet number: 20
\scanpage{22}
\mspage{3}

\[
\mathrm{Ob}\,X\simeq F=M/\unsure{\mathcal R|R}.
\]

Si $x,y\in\mathrm{Ob}\,X$, on a
\[
\mathrm{Hom}(x,y)\simeq\{u\in M\mid x=yu\},
\]
la composition des flèches \deleted{\unsure{étant}} donnée par la composition de $M$. La catégorie $X$ admet $a=\dot 1$ comme plus grand élément, tout autre objet $x$ étant de la forme
\[
x=\dot u=au.
\]
\inserted{Notons que $Y$ étant pseudo-cofiltrant, il en est de $=$ de $\deleted{\unsure{Y/F}}X=Y_{/F}$ (en fait, pour \unsure{tt} $F\in\mathrm{Ob}\,Y^\wedge$).}

À quelles conditions $X$ admet-il un plus petit élément $e$, lequel sera de la forme
\[
e=a\psi\qquad(\psi\in M\ \text{fixé})\,?
\]

On voit tout de suite que la condition est la suivante
\[
\tag*{\textcircled{A}}\begin{aligned}&\forall u\in M,\ \exists\,u'\in M\quad\text{tel que}\quad uu'\overset{\mathcal R}\sim\psi,\\ &\bigl(\text{i.e }\exists\,p,q\in\mathbb N\ \text{avec}\quad \varphi^p\psi=\varphi^q uu'\bigr).\end{aligned}
\]

Par exemple, il suffit que $\psi$ soit un élément \uline{minimal} de $M$ (i.e. $\forall u,\exists u'$ avec $\psi=uu'$).

% archive sheet number: 21
\scanpage{23}
\mspage{4}

Posons
\[
\begin{aligned}M_0=\mathrm{End}_X(e)=\{\,f\in M\mid ef&=\deleted{\unsure{f}}e\ \text{i.e.}\ a\psi f=a\psi\\ &\text{i.e.}\ \unsure{a\psi|\psi}f\overset{R}\sim\psi\,\}\end{aligned}
\]

Ainsi, on a pour $f\in M$
\[
f\in M_0\quad\Longleftrightarrow\quad\forall\,p,q\in\mathbb N\qquad\varphi^p\psi f=\varphi^q\psi.
\]

Je vais écrire une condition sur $\varphi,\psi$ qui impliquera que pour $f\in M_0$ fixé, l'entier $\delta(f)\overset{\mathrm{def}}{=}q-p$ ne dépend pas du choix de $p,q$.

Supposons donc qu'on ait
\[
\left\{\begin{aligned}\varphi^p\psi f&=\varphi^q\psi\\ \varphi^{p'}\psi f&=\varphi^{q'}\psi\end{aligned}\right.
\]
on en conclut
\[
\varphi^{p+p'}\psi f=\left\{\begin{aligned}&\varphi^{p'+q}\psi\\ &\varphi^{p+q'}\psi\end{aligned}\right.
\]
donc
% letterhead
\par\smallskip
{\footnotesize\noindent\begin{minipage}[t]{.46\linewidth}\raggedright
MINISTERE DE L'EDUCATION NATIONALE\\ ACADEMIE D'AIX-MARSEILLE\\ \rule{2.6em}{.3pt}\\ Inspection Académique\\ de Vaucluse\\ \rule{2.6em}{.3pt}\\ Tél. : 90 82 99 76\\ Référence : Secrétariat n\textsuperscript{o}
\end{minipage}\hfill%
\begin{minipage}[t]{.51\linewidth}\raggedright
AVIGNON, le\\ L'Inspecteur d'académie,\\ directeur des services départementaux\\ de l'Education Nationale de Vaucluse,\\ à M
\end{minipage}\par}

% archive sheet number: 22
\scanpage{24}
\mspage{5}

\[
\varphi^{p'+q}\psi=\varphi^{p+q'}\psi
\]
\margin{On veut conclure $p'+q=p+q'$ i.e. $q-p=q'-p$}

On va supposer
\textcircled{B} Si $r,s\in\mathbb N$ sont tels que $\varphi^r\psi=\varphi^s\psi$, on a $r=s$.

On trouve alors une application
\diagram{\begin{tikzcd}[ampersand replacement=\&]M_0\arrow[r,"\delta"]\&\mathbb Z\quad\deleted{\unsure{g}},\end{tikzcd}}
et il est immédiat que celle-ci est \unsure{multiplicative|multiplication} ; i.e. $\delta(fg)=\delta(f)+\delta(g)$ :
\[
\left\{\begin{aligned}\varphi^p\psi f&=\varphi^q\psi\\ \varphi^{p'}\psi g&=\varphi^{q'}\psi\end{aligned}\right.
\]
\margin{multiplier par $g$ à droite}
implique $\varphi^p\psi fg=\varphi^q\psi g$
et en multipliant par $\varphi^{p'}$ à gauche
\[
\varphi^{p+p'}\psi fg=\varphi^{q+p'}\psi g=\varphi^q\bigl(\underbrace{\varphi^{p'}\psi g}_{\varphi^{q'}\psi}\bigr)=\varphi^{q+q'}\psi
\]

donc $\delta(fg)=(q+q')-(p+p')=\delta(f)\unsure{+|-}\delta(g)$.

Donc $\delta$ se factorise en
\diagram{\begin{tikzcd}[ampersand replacement=\&,column sep=large]M_0\arrow[r,"\mathrm{can}"]\arrow[rr,bend right=43,"\delta"']\&\underset{\substack{\text{groupe}\\\text{enveloppant}\\\text{de }M_0}}{G_0}\arrow[r,"\delta_0"]\&\mathbb Z.\end{tikzcd}}

% archive sheet number: 23
\scanpage{25}
\mspage{6}

Je vais supposer
\textcircled{C} L'hom \diagram{\begin{tikzcd}[ampersand replacement=\&]M_0\arrow[r,no head,"\delta"]\&\mathbb Z\end{tikzcd}} n'est pas nul,
\[
\left|\begin{aligned}&\text{i.e }\exists\,f\in M_0\quad\text{et}\quad p,q\in\mathbb N,\quad\text{avec}\\ &\varphi^p\psi f=\varphi^q\psi\quad\text{et}\quad p\ne q\end{aligned}\right.
\]

\uline{Exemple.} Si $\psi\sim\deleted{\unsure{f}}1$, alors \deleted{$\unsure{\varphi f\sim1}$ \unsure{sera}} $\underset{\substack{\unsure{?|2}\\ f}}{\psi f}\sim\psi$ \deleted{et la relation $\underset{\unsure{?}}{\psi}\underset{f}{f}\sim\underset{\substack{?\\1}}{\psi}$} équivaut à $f\sim1$ \inserted{i.e $f\in M_0$}. \deleted{Si on prend $f=\varphi$,} On vérifie aisément que $\delta(f)$ est \unsure{donné|déterminé}, \deleted{\unsure{aura} $\underbrace{\psi}_{1}\underbrace{f}_{\varphi}=\varphi$ \unsure{de} $\psi f=\varphi\psi$ i.e} en termes de $\varphi^p f=\varphi^q$, par $\delta(f)=q-p$, \deleted{\unsure{avec} $\unsure{\psi f=\varphi\psi}$} en particulier \deleted{$\delta(f)=\delta(\varphi)=1$,} on aura $\delta(\varphi)=1$ donc $\delta$ est surjectif. \deleted{donc $\delta$ est surjectif, et la condition \textcircled{C} est vérifiée.}

\uline{NB} Dans le cas, $a\in\mathrm{Ob}\,X$ est à la fois objet minimal et maximal.

La condition \textcircled{C} implique que $G_0=\pi_1(X_0,e)$ admet un quotient isomorphe à $\mathbb Z$ (savoir l'image de $G_0$ dans $\mathbb Z$ par $\delta_0$).

\[
\underline{\varphi^\alpha\psi f}=\varphi^\beta f
\]
\[
\psi\varphi\sim\psi
\]
\[
\varphi^\alpha\psi\sim\varphi^\beta
\]
% letterhead
\par\smallskip
{\footnotesize\noindent\begin{minipage}[t]{.46\linewidth}\raggedright
\deleted{MINISTERE DE L'EDUCATION NATIONALE}\\ \deleted{ACADEMIE D'AIX-MARSEILLE}\\ \rule{2.6em}{.3pt}\\ \deleted{Inspection Académique}\\ \deleted{de Vaucluse}\\ \rule{2.6em}{.3pt}\\ \deleted{Tél. : 90 82 99 76}\\ \deleted{Référence : Secrétariat n\textsuperscript{o}}
\end{minipage}\hfill%
\begin{minipage}[t]{.51\linewidth}\raggedright
\deleted{AVIGNON, le}\\ \deleted{L'Inspecteur d'académie,}\\ \deleted{directeur des services départementaux}\\ \deleted{de l'Education Nationale de Vaucluse,}\\ \deleted{à M}
\end{minipage}\par}

% archive sheet number: 24
\scanpage{26}
\mspage{7}

On va maintenant appliquer la prop. 1$^{*}$, pour calculer $\pi_1(X,e)$, en termes de $\pi_1(X_0,e)=G_0$, et des relations explicitées par ailleurs. Je veux \deleted{\unsure{prouver}} \inserted{montrer} que les images de ces relations dans $\mathbb Z$ par $\delta_0$ sont nulles, donc que $\pi_1(X_0,e)$ admet un groupe quotient isomorphe à $\mathbb Z$, donc n'est pas $\sim\{1\}$, donc que $X_0$ n'est pas $1$-connexe.
% slanted
\margin{Si on applique le NB page 2, \deleted{\illegible{1}} les développements jusqu'à la page 12 sont inutiles !}

Je reprends donc le diagramme
\diagram{\begin{tikzcd}[ampersand replacement=\&] \& x\arrow[r,"f"]\&y\arrow[r,"g"]\&z\\ \text{(*)}\&\&e\arrow[ul,"u"]\arrow[u,"v"]\arrow[ur,"w"']\&{,}\end{tikzcd}}

en posant
\[
\left\{\begin{aligned}z&=a\rho,\quad\text{donc}\\ y&=a\rho g\\ x&=a\rho gf\end{aligned}\right.
\]
\[
\left\{\begin{aligned}e&=xu=a\rho gfu\\ e&=yv=a\rho gv\\ e&=zw=a\rho w\end{aligned}\right.
\]
et comme $e=a\psi$, les trois relations précédentes s'écrivent
\[
\tag*{(**)}\boxed{\rho gfu\sim\rho gv\sim\rho w\sim\psi}
\]

La donnée d'un diagramme \unsure{cas|($*$)} dans $X$, avec $z=a\rho$ (pour $\rho\in M$ fixé) équivaut à la donnée des six éléments

% archive sheet number: 25
\scanpage{27}
\mspage{8}

\[
\rho,f,g,u,v,w\qquad\in M
\]
de $M$, satisfaisant les relations ($**$), ou encore tels qu'il existe $p,q,r,s\in\mathbb N$ avec
\[
\boxed{\varphi^p(\rho gfu)=\varphi^q(\rho gv)=\varphi^r(\rho w)=\varphi^s\psi}
\]

Ceci posé, on doit poser la relation
\[
\boxed{[w,gv][v,fu][w,gfu]^{-1}=1},
\]

Pour l'expliciter, considérons de façon générale
\diagram{\begin{tikzcd}[ampersand replacement=\&]e\arrow[r,shift left=.55ex,"\lambda"]\arrow[r,shift right=.55ex,"\mu"']\&\xi=a\sigma\end{tikzcd}}
donc $e=a\sigma\lambda=a\sigma\mu$ i.e.
\[
\boxed{\sigma\lambda\sim\sigma\mu\sim\psi},
\]

et choisissons $\lambda',\mu'\in M$ tels que $\lambda\lambda'=\mu\mu'=\unsure{\pi|\beta}$ de sorte qu'on a un diagramme commutatif
\diagram{\begin{tikzcd}[ampersand replacement=\&]\&e\arrow[dr,"\lambda"]\&\\ \eta\arrow[ur,"\lambda'"]\arrow[dr,"\mu'"']\&\&\xi\\ \&e\arrow[ur,"\mu"']\&\end{tikzcd}}
\[
\begin{aligned}\text{avec}\quad\eta&=e\lambda'=\xi\lambda\lambda'=a\sigma\lambda\lambda'\deleted{\unsure{=\pi}}\inserted{=a\sigma\pi}\\ \cdot&=e\mu'=\xi\mu\mu'=a\sigma\mu\mu'=a\sigma\pi\end{aligned}
\]

on va \deleted{choisir} prendre une flèche \diagram{\begin{tikzcd}[ampersand replacement=\&]e\arrow[r,"\nu"]\&\eta\end{tikzcd}} donc un $\nu\in M$ tel que
% letterhead
\par\smallskip
{\footnotesize\noindent\begin{minipage}[t]{.46\linewidth}\raggedright
MINISTERE DE L'EDUCATION NATIONALE\\ ACADEMIE D'AIX-MARSEILLE\\ \rule{2.6em}{.3pt}\\ Inspection Académique\\ de Vaucluse\\ \rule{2.6em}{.3pt}\\ Tél. : 90 82 99 76\\ Référence : Secrétariat n\textsuperscript{o}
\end{minipage}\hfill%
\begin{minipage}[t]{.51\linewidth}\raggedright
AVIGNON, le\\ L'Inspecteur d'académie,\\ directeur des services départementaux\\ de l'Education Nationale de Vaucluse,\\ \deleted{à M}
\end{minipage}\par}

% archive sheet number: 26
\scanpage{28}
\mspage{\textcircled{9}}

\[
e=\eta\nu=\deleted{\unsure{a\sigma|\text{avec}}}a\sigma\pi\nu\qquad\text{\unsure{i.e.|i.k.} tel que}
\]
\[
\boxed{\sigma\pi\nu\deleted{\unsure{=}}\sim\psi}\qquad\text{où }\pi=uu'=vv'.
\]

Alors on aura
\[
\left\{\begin{aligned}&\lambda'\nu,\deleted{\unsure{\lambda|\&}}\mu'\nu\in M_0\qquad\text{et}\\ &\lambda(\lambda'\nu)=\mu(\mu'\nu)\end{aligned}\right.
\]
d'où
\[
[\mu,\lambda]=(\mu'\nu)(\lambda'\nu)^{-1}.
\]

L'existence de $\nu\deleted{\unsure{{}_{*}|{,}}}$ satisfaisant $(\sigma\pi)\nu\sim\psi$ est garantie par \textcircled{A}, on choisit (pour tout $u\in M$) un $\unsure{\Theta|\theta}(u)$ tel que $u\Theta(u)\sim\psi$, et on prendra ici $\nu=\Theta(\sigma\pi)$. Donc on aura
\[
\left\{\begin{aligned}&[\mu,\lambda]=(\mu'\Theta(\sigma\pi))(\lambda'\Theta(\sigma\pi))^{-1}\\ &\text{où }\lambda',\mu'\in M\ \text{sont tels que}\quad\lambda\lambda'=\mu\mu'(\overset{\mathrm{def}}{=}\pi)\\ &\text{et }\sigma\inserted{\in M}\ \text{est tel que}\quad\xi=a\sigma.\end{aligned}\right.
\]

Je crois que je vais craquer, et restreindre la généralité, en supposant \inserted{carrément} que $\psi$ est un objet minimal\unsure{,|e} (ce qui implique A), donc que pour tt $u$ existe $\Theta(u)$ tel que $u\Theta(u)=\psi$ (pas seulement $\sim\psi$) donc ci-dessus on peut prendre $\lambda'=\Theta(\lambda)$, $\mu'=\Theta(\mu)$, et on aura $\lambda\lambda'=\mu\mu'\overset{\mathrm{def}}{=}\pi=\psi$.

% archive sheet number: 27
\scanpage{29}
\mspage{10}

\[
[\mu,\lambda]=\bigl(\underbrace{\unsure{\Theta|\theta}(\mu)\Theta\sigma\psi)}_{\in M_0}\bigr)\bigl(\underbrace{\Theta(\lambda)\Theta(\sigma\psi)}_{\in M_0}\bigr)^{-1}
\]

On applique ceci aux quatre crochets qui figurent \unsure{dans|sur} la relation de la page 8 :
\[
\begin{aligned}[w,gv]&=\bigl(\Theta(w)\Theta(\rho\psi)\bigr)\bigl(\Theta(gv)\Theta(\rho\psi)\bigr)^{-1}\\ [v,fu]&=\bigl(\Theta(v)\Theta(\rho g\psi)\bigr)\bigl(\Theta(fu)\Theta(\rho g\psi)\bigr)^{-1}\\ [w,gfu]&=\bigl(\Theta(w)\Theta(\rho\psi)\bigr)\bigl(\Theta(gfu)\Theta(\rho\psi)\bigr)^{-1}\end{aligned}
\]

La relation devient dès lors
\[
\boxed{\begin{aligned}&\bigl(\Theta(gv)\Theta(\rho\psi)\bigr)^{-1}\bigl(\Theta(v)\Theta(\rho g\psi)\bigr)\bigl(\Theta(fu)\Theta(\rho g\psi)\bigr)^{-1}\\ &\hspace{10em}\inserted{\Theta(gfu)^{-1}\Theta(\rho\psi)}=1\end{aligned}}
\]

(\uline{NB} $w$ y a disparu sans laisser de \unsure{trace|traces}) En fait, on a besoin d'un peu plus, pour arriver à s'en sortir :
% letterhead
\par\smallskip
{\footnotesize\noindent\begin{minipage}[t]{.46\linewidth}\raggedright
\deleted{MINISTERE DE L'EDUCATION NATIONALE}\\ \deleted{ACADEMIE D'AIX-MARSEILLE}\\ \rule{2.6em}{.3pt}\\ \deleted{Inspection Académique}\\ \deleted{de Vaucluse}\\ \rule{2.6em}{.3pt}\\ \deleted{Tél. : 90 82 99 76}\\ \deleted{Référence : Secrétariat n\textsuperscript{o}}
\end{minipage}\hfill%
\begin{minipage}[t]{.51\linewidth}\raggedright
\deleted{AVIGNON, le}\\ \deleted{L'Inspecteur d'académie,}\\ \deleted{directeur des services départementaux}\\ \deleted{de l'Education Nationale de Vaucluse,}\\ \deleted{à M}
\end{minipage}\par}

% archive sheet number: 28
\scanpage{30}
\mspage{11}

On va supposer
\[
\boxed{\psi\sim1\qquad\text{i.e}\quad\exists\,\alpha,\beta\in\mathbb N}\quad\text{avec}\quad\boxed{\varphi^\alpha\psi=\varphi^\beta}
\]
% slanted
\margin{ce qui implique \unsure{C|\textcircled{C}|\textcircled{D}}}

Donc les relations
\[
\rho gfu\sim\rho gv\sim\rho w\sim\psi
\]
équivalent à
\[
\rho gfu\sim\rho gv\sim\rho w\sim1
\]
i.e. à
\[
\boxed{\rho gfu,\rho gv,\rho w\in M_0.}
\]

(Que $w$ ait disparu de la relation \deleted{\unsure{fin}} \unsure{ci-dessus|fondamentale} ne doit pas nous inquiéter, c'est une relation en $\deleted{\unsure{\rho,v,g,f,|f,\psi,g,q}}\underline{\rho,f,g,u,v}$, et où on peut poser $w=gv$, de sorte qu'il \unsure{en|ne} reste, comme \unsure{conditions|condition} sur ces cinq variables en question, que
\[
\underbrace{\rho gfu}_{\substack{\mathrm{def}\,\Vert\\\pi}},\ \underbrace{\rho gv}_{\substack{\Vert\,\mathrm{def}\\\pi'}}\in M_0,
\]
donc on trouve deux \unsure{élmts|éléments} $\pi,\pi'\in M_0$.

Ceci posé, on trouve les relations
\unsure{$\cdot\vert$|i)}

% archive sheet number: 29
\scanpage{31}
\mspage{12}

\[
\left\{\begin{aligned}\underbrace{\pi}_{\substack{\Vert\\\rho gfu}}\bigl(\underbrace{\unsure{\Theta|\theta}(gv)\Theta(\rho\psi)}_{c_1}\bigr)&=\pi\bigl(\underbrace{\Theta(v)\Theta(\rho g\psi)}_{c_2}\bigr)=\psi\\ \pi'\bigl(\underbrace{\Theta(fu)\Theta(\rho g\psi)}_{c_3}\bigr)&=\pi'\bigl(\underbrace{\Theta(gfu)\Theta(g\psi)}_{c_4}\bigr)=\psi\end{aligned}\right.
\]

d'où
\[
\begin{aligned}\delta(c_1)=\delta(c_2)&=\delta(\psi)-\delta(\pi)\\ \delta(c_3)=\delta(c_4)&=\delta(\psi)-\delta(\pi')\end{aligned}
\]

d'où
\[
\begin{aligned}\delta_0(c_1c_2^{-1}c_3c_4^{-1})&=\bigl(\underbrace{\delta(c_1)-\delta(c_2)}_{\substack{\Vert\\0}}\bigr)+\bigl(\underbrace{\delta(c_3)-\delta(c_4)}_{\substack{\Vert\\0}}\bigr)\\ &=0,\end{aligned}
\]

donc $\deleted{G_0\overset{\delta}{\text{---}}}$ l'hom surjectif $\delta_0:G_0\longrightarrow\mathbb Z$ se factorise en un hom surjectif
\[
\delta'_0:\pi_1(X,e)\longrightarrow\mathbb Z.
\]
% letterhead
\par\smallskip
{\footnotesize\noindent\begin{minipage}[t]{.46\linewidth}\raggedright
MINISTERE DE L'EDUCATION NATIONALE\\ ACADEMIE D'AIX-MARSEILLE\\ \rule{2.6em}{.3pt}\\ Inspection Académique\\ de Vaucluse\\ \rule{2.6em}{.3pt}\\ Tél. : 90 82 99 76\\ Référence : Secrétariat n\textsuperscript{o}
\end{minipage}\hfill%
\begin{minipage}[t]{.51\linewidth}\raggedright
AVIGNON, le\\ L'Inspecteur d'académie,\\ directeur des services départementaux\\ de l'Education Nationale de Vaucluse,\\ à M
\end{minipage}\par}

% archive sheet number: 30
\scanpage{32}
\mspage{13}

En résumé :

\uline{Proposition 2} : \unsure{Soient|Soit} $M$ un monoïde, \deleted{On suppose} $\varphi$ et $\psi$ deux éléments de $M$.

On suppose :

a) $\psi$ est un élément minimal de $M$, i.e pour \deleted{tout} $u\in M$, existe $\unsure{\Theta|\theta}(u)\in M$ avec $\psi=u\Theta(u)$. (Donc $M$ i.e. $Y=B_M$ est pseudo-cofiltrant, i.e. $Y^\wedge$ est loc. irréductible.)

b) $\psi\sim1$ pour la relation d'équivalence invariante à droite $R_\varphi$ définie par $\varphi$
\[
(u\overset{R_\varphi}{\sim}v\quad\Longleftrightarrow\quad\exists\,p,q\in\mathbb N\ \text{avec}\ \varphi^p u=\varphi^q v),
\]
i.e $\exists\,\alpha,\beta\in\mathbb N$ avec $\varphi^\alpha\psi=\varphi^\beta\deleted{\unsure{\psi}}$.

c) Si $\deleted{p,q}\unsure{r,s|p,q}\in\mathbb N$ sont tels que $\varphi^{\deleted{\unsure{p}}r}\psi=\varphi^s\psi$, alors $r=s$.

Soit $Y=B_M$, $F=M/\unsure{\mathcal R_\varphi|R_\varphi}$ considéré comme $M$-ensemble à droite i.e. comme objet de $Y^\wedge$, $X=Y_{/F}$. Alors l'objet $e=\dot 1$ de $X$ est à la fois un plus \deleted{petit}\inserted{grand} et un plus petit élément de $X$.

De plus, $\deleted{\pi_{1}(X,e)}$ il existe un hom surjectif $\pi_1(X,e)\longrightarrow\mathbb Z$, en particulier $X$ n'est pas $1$-connexe.

% archive sheet number: 31
\scanpage{33}
\mspage{14}

\uline{Corollaire} Supposons que $M$ soit \deleted{de la} isomorphe, soit au monoïde $\mathrm{Ep}(E)$ où $E$ est un ens infini, soit au monoïde opposé à $\mathrm{Mon}(E)$, avec $E$ un ens. infini. Alors $M$ est pseudocofilt\deleted{\unsure{rant}}\inserted{(il admet $=$ des éléments minimaux)} mais non cofiltrant (en fait $\unsure{f|g}uf=vf\Longrightarrow u=v$), $Y=B_M$ est $1$-connexe, mais il existe $F\in\mathrm{Ob}\,Y^\wedge$ tel que $X=Y_{/F}$ soit $0$-connexe et non $1$-connexe.

Pour ceci, il suffit de prendre \deleted{pour} la construction de la proposition, en prenant \deleted{\unsure{pour}} $\psi=\varphi$ n'importe quel objet minimal de $M$.

Il suffit de vérifier qu'on ne peut avoir $\varphi^{\deleted{\unsure{p}}r}=\varphi^s$ avec $r\ne s$. Mais si $r=s+h$, $h>0$, on aurait $\varphi^h\varphi^s=\mathrm{id}\,\varphi^s$, d'où $\varphi^h=1$, donc $\varphi$ serait inversible.
% letterhead
\par\smallskip
{\footnotesize\noindent\begin{minipage}[t]{.46\linewidth}\raggedright
\deleted{MINISTERE DE L'EDUCATION NATIONALE}\\ \deleted{ACADEMIE D'AIX-MARSEILLE}\\ \rule{2.6em}{.3pt}\\ \deleted{Inspection Académique}\\ \deleted{de Vaucluse}\\ \rule{2.6em}{.3pt}\\ \deleted{Tél. : 90 82 99 76}\\ \deleted{Référence : Secrétariat n\textsuperscript{o}}
\end{minipage}\hfill%
\begin{minipage}[t]{.51\linewidth}\raggedright
\deleted{AVIGNON, le}\\ \deleted{L'Inspecteur d'académie,}\\ \deleted{directeur des services départementaux}\\ \deleted{de l'Education Nationale de Vaucluse,}\\ \deleted{à M}
\end{minipage}\par}

% archive sheet number: 32
\scanpage{34}
\mspage{15}

Mais un élément inversible ne peut être minimal. En effet, on aurait pour tt $u\in M$ un $u'\in M$ tel que $\varphi=uu'$, \deleted{i.e.\ $\unsure{u}$} \unsure{en|donc|et en} posant $v=u'\varphi^{-1}$, on aurait
\[
uv=1
\]
montrons que l'on aurait aussi
\[
p=vu=1,
\]
en effet on a $p^2=p$, i.e. $p.p=\mathrm{id}.p$ donc $p=\mathrm{id}$. Donc tous les objets de $M$ seraient inversibles, i.e. $M$ serait un groupe, ce qui n'est pas le cas.

\uline{NB} Considérons \deleted{une partie} \inserted{(sous les conditions du corollaire) un} sous-$M$-ensemble à droite $F$ de $M$, donc $F$ \deleted{s'identifie}\inserted{peut se} regarder comme un objet de $Y^\wedge$ ($Y=B_M$), d'où une catégorie $X=Y_{/F}$. On peut montrer que les $X$ sont tous $1$-connexes.

J'ai eu du mal à construire un $F$ \deleted{tel que}\inserted{$0$-connexe} qui ne soit $1$-connexe !

% archive sheet number: 33
\scanpage{35}
% letterhead
\par\smallskip
{\footnotesize\noindent\begin{minipage}[t]{.46\linewidth}\raggedright
MINISTERE DE L'EDUCATION NATIONALE\\ ACADEMIE D'AIX-MARSEILLE\\ \rule{2.6em}{.3pt}\\ Inspection Académique\\ de Vaucluse\\ \rule{2.6em}{.3pt}\\ Tél. : 90 82 99 76\\ Référence : Secrétariat n\textsuperscript{o}
\end{minipage}\hfill%
\begin{minipage}[t]{.51\linewidth}\raggedright
AVIGNON, le\\ L'Inspecteur d'académie,\\ directeur des services départementaux\\ de l'Education Nationale de Vaucluse,\\ à M
\end{minipage}\par}

% archive sheet number: 34
